\section{The conductor-sixteen sectors in the actual growing arithmetic window}
\label{sec:conductor16-actual-window-reduction}

\paragraph{Status.}
This proof applies the original strict-zero theorem AC-G142 to the
actual arithmetic window.  Its class certificate is the separately proved
AC-G142-class; no provisional class group is used.  The remaining
arithmetic bound and mixed theta realization are still unproved.
All other hypotheses in AC-K35 remain unchanged.  No new coefficient-prime
uniformity is inferred from a calculation at $269$.

Put $p=269$, $K=\mathbb Q(\omega)$ and $t^3=p$.
Let $\chi:G_K\to\mu_3$ be the character of $K(t)/K$ fixed throughout
the cubic argument.  Write
\[
 F_{16}=\mathbb Q(w),\qquad w=\zeta_{16}-\zeta_{16}^{-1},
 \qquad w^4+4w^2+2=0.
 \tag{CW.1}
\]
Choose the odd quartic Dirichlet character $\phi$ of conductor $16$
whose kernel in $(\mathbb Z/16)^\times$ is $\{1,7\}$ and for which
$\phi(13)=i$.  In particular $\phi(5)=-i$ and $\phi(-1)=-1$.
Either choice of the fourth root $i$ gives the same pair
$\{\phi,\phi^{-1}\}$.  We use arithmetic reciprocity for the finite
Galois character and the actual cyclotomic character $c_p$.

\paragraph{The exact strict-zero theorem being applied.}
AC-G142, with $r=0$ and $\kappa=1$, supplies the following statement,
denoted $Z_{16}$:
for every finite coefficient discrete valuation ring $O$ of residual
characteristic $p$ containing the values of $\chi$ and $\phi$, with
fraction field $E$, and every
\[
 j\in\{1,2\},\quad \sigma\in\{1,-1\},\quad
 e\geq1,\quad e\equiv267\pmod{268},
 \tag{CW.2}
\]
put $A=(E/O)(\chi^{-j}\phi^{-\sigma}c_p^{-e})$.  The required
vanishing and its local condition are
\[
 \begin{gathered}
 \mathcal S_{\Sigma_0}(A)=0,\qquad
 \mathcal S_\Sigma(A):=
 \ker\bigl(H^1(G_{K,\Sigma},A)\longrightarrow\mathcal L_\Sigma(A)\bigr),\\
 \mathcal L_\Sigma(A):=
 \bigoplus_{v\mid p}H^1(K_v,A)
       \oplus\!\bigoplus_{\substack{v\in\Sigma_{\rm fin}\\v\nmid p}}
                                      H^1(I_v,A).
 \end{gathered}
 \tag{CW.3}
\]
where $\Sigma_0$ contains the complex place and the places over $2,3,p$.
It is enough to prove this for that minimal displayed set; the proof
below checks extension to every actual set of primes used by AC-K25.
The condition at $p$ is the full local condition, not merely vanishing
on inertia or an ordinary quotient.  There is no automorphic seed,
native support ideal, original-model hypothesis or mixed theta input
in $Z_{16}$.

Use the actual fields and pages of (FG.3) and AC-K35:
\[
 \begin{split}
 F_u&=\mathbb Q(\zeta_{2^{h(u)+2}},\zeta_7,
                   \sqrt{-q},\sqrt{-d},\omega),\\
 H_u&=F_u(t),\qquad G_u=\operatorname{Gal}(H_u/K),\\
 u^{2+8\varepsilon}&\leq2^{h(u)}\leq2u^{2+8\varepsilon},
 \qquad 0<\varepsilon<1/256.
 \end{split}
 \tag{CW.4}
\]
The pair $q,d$ is the one prescribed by the source window.  No new
auxiliary prime or new field is chosen in this argument.
Let $D_u=|G_u|$, let $O_\ell,b_\ell$ be exactly the coefficient rings
and degrees in (FC.1), and put $T(u)=\lfloor u^{399}\rfloor$.
For $u\geq2$, define the subset of the actual index labels
\[
 \mathscr Z_u=
 \{(p,\chi^{-j}\phi^{-\sigma},e):
       j=1,2,\ \sigma=\pm1,\ 1\leq e<T(u),\
       e\equiv267\pmod{268}\}.
 \tag{CW.5}
\]
In this formula the finite characters are inflated to $G_u$.
Let $I_{399}^{\circ}(H_u,u)$ be the sum (TC.1), with precisely the
labels in $\mathscr Z_u$ omitted.  Every other character, every other
weight and every other coefficient prime is retained.

\begin{theorem}[AC-G143-window: removing proved sectors from the actual window]
\label{AC-G143-window}
The following statements hold on every actual host (CW.4), uniformly
in its allowed page and pair $q,d$.

First, every label in $\mathscr Z_u$ is an actual character label of
AC-K25, and its strict group (FC.2) is zero.  There are exactly
\[
                         4\left\lfloor\frac{T(u)}{268}\right\rfloor
 \tag{CW.6}
\]
such labels.  Their total contribution to the strict sum (FC.15) is
exactly zero.  Their contribution to the derived index sum need not
be zero, but it satisfies
\[
 \begin{split}
 0\leq I_{399}(H_u,u)-I_{399}^{\circ}(H_u,u)
 \leq{}&8D_u^2T(u)^2\log(8T(u)D_u)\\
 &+\tfrac12D_u\mathcal R_u T(u)^2
       +2D_u r_u c_0 T(u),
 \end{split}
 \tag{CW.7}
\]
where $\mathcal R_u,r_u,c_0$ are the exact quantities in (FL.1).
In particular this difference is
\[
 O_\varepsilon\bigl(u^{802+16\varepsilon}(\log u)^2\bigr)
       =o\bigl(\exp(u^\varepsilon/2)\bigr).
 \tag{CW.8}
\]

Second, the numerical premise (TC.3) may be replaced by the identical
premise with $I_{399}^{\circ}$ in place of $I_{399}$, with positive
relative slack and every common-window quantifier unchanged.
The old premise implies the new one with the same parameters.
The new premise implies the old one with half its slack after
increasing only the initial threshold.  Consequently it suffices
for the same conditional conclusion of AC-K35, retaining the same
unproved $\mathcal R_{\rm cub}$ and every other source hypothesis.
This theorem does not assert an upper bound for the remaining sum.
\end{theorem}

\begin{proof}
\proofstep{1. The characters really occur on the prescribed host.}
The formula for $w$ in (CW.1) follows from
$w^2=-2+\sqrt2$.  The automorphism $\zeta_{16}\mapsto\zeta_{16}^{13}$
sends $w$ to $-w^3-3w$ and has square $w\mapsto-w$.
The subgroup of $(\mathbb Z/16)^\times$ fixing $w$ is $\{1,7\}$.
Thus $F_{16}$ is exactly the cyclic quartic field cut out by $\phi$;
complex conjugation acts as the square of its chosen generator.

If $u\geq2$, then the page inequality gives $2^{h(u)}\geq4$,
so $h(u)\geq2$ and $\mathbb Q(\zeta_{16})\subset F_u$.
Also $K(t)\subset H_u$ by definition.  Therefore $\chi$ and
$\phi|_{G_K}$ factor through $G_u$.  The intersection
$K(t)\cap F_u=K$ is the actual disjointness proved in AC-K25:
otherwise the nonabelian $S_3$ splitting field $K(t)/\mathbb Q$
would be contained in the abelian extension $F_u/\mathbb Q$.
The four characters in (CW.5) are distinct: their restrictions
to $\operatorname{Gal}(K(t)/K)$ distinguish $j$, and the quotient
through $F_{16}$ distinguishes $\sigma$.

The exponent of $G_u$ is $N=3\cdot2^{h(u)}$, so $O_p$ contains
both the cubic and quartic values.  Since $p\nmid N$, this ring is
unramified at $p$, but no assertion of degree one or a fixed residue
degree is needed.  Its actual rank $b_p$ is retained throughout.
The coefficient in (FC.2) is precisely
\[
 (O_p\otimes_{\mathbb Z_p}\mathbb Q_p/\mathbb Z_p)
                       (\chi^{-j}\phi^{-\sigma})(-e)
     =(E_p/O_p)(\chi^{-j}\phi^{-\sigma}c_p^{-e}),
 \tag{CW.9}
\]
where $E_p=\operatorname{Frac}(O_p)$.  The tensor identification
is the one induced by the finite free $\mathbb Z_p$-module $O_p$.
Thus no change of the original coefficient lattice or implicit
division by its residue degree occurs.

\proofstep{2. The local conditions agree, including added primes.}
Let $A$ be the coefficient in (CW.9), let $\Sigma$ be the actual
set $\Sigma_p$ from AC-K25, and enlarge it if necessary by the
fixed finite set $\Sigma_0$.  This enlargement does not alter the
strict kernel.  Here is the elementary comparison that also handles
the primes over $7,q,d$ already in $\Sigma$.

If $v$ is an added prime, $v\nmid p$ and $A$ is unramified there.
For a cocycle representing a strict class, its restriction to $I_v$
has zero cohomology class.  Since $I_v$ acts trivially on $A$,
every coboundary on $I_v$ is zero, so this restriction is the zero
cocycle.  The closed normal subgroup generated by the added inertia
groups acts trivially on $A$.  The cocycle identity then gives zero
on every conjugate and on their closed normal closure.  Consequently
the class descends uniquely to the Galois group with these primes
removed.  Conversely inflation preserves every displayed local
condition.  Inflation is injective by the first term of the
inflation--restriction sequence.  This proves equality of the
strict kernels before and after enlargement.

The coefficient's possible ramification is contained in $2,3,p$;
the cyclotomic twist is unramified away $p$.  At $p$ both definitions
require zero in the full $H^1(K_v,A)$.  At every other prime both
require zero on inertia.  Hence $Z_{16}$ gives
$\mathcal S_{p,\chi^{-j}\phi^{-\sigma}}(e)=0$ on the actual host.
Writing $e=268n-1$, the condition $1\leq e<T(u)$ is equivalent to
$1\leq n\leq\lfloor T(u)/268\rfloor$.  This proves (CW.6).

\proofstep{3. The exact positive-index comparison.}
For one removed label set
$C=R\Gamma(G_{K,\Sigma_p},O_p(\eta^{-1})(e+1))$.
The dual lattice is the actual codifferent dual in AC-K26-local,
identified with this free rank-one $O_p$-lattice as in (FL.2).
The derived determinant statement (EU.2) gives, term by term,
\[
 \frac{\log\#k_{O_p}}{b_p}\max(0,i_{p,\eta,e})
       \leq \frac1{b_p}\log\#H^2(C).
 \tag{CW.10}
\]
Because the strict group at this label has order one, the last
quantity is exactly
\[
 \frac1{b_p}\left|\log\#H^2(C)
                   -\log\#\mathcal S_{p,\eta}(e)\right|.
 \tag{CW.11}
\]
Sum over the four removed characters and over their indicated
weights.  All terms in (FL.4) are nonnegative, so this partial sum
is bounded by the sum of (FL.4) over all $1\leq e<T(u)$.
The inequalities
\[
 \sum_{1\leq e<T}e\leq T^2/2,\quad
 \log(8eD)\leq\log(8TD),\quad T-1\leq T
\]
give exactly (CW.7).  No claim that strict vanishing forces a
derived index to vanish has been used.  In particular no termwise
identification of $H^2$ with the strict group, and no omission of
the local comparison terms, is being made.

The actual host estimates from (FL.21), (FG.11) and the source
window inequalities are
\[
 D_u\leq288u^{2+8\varepsilon},\quad r_u\leq12,\quad
 \mathcal R_u=O_\varepsilon(u^{\varepsilon/5}(\log u)^2).
 \tag{CW.12}
\]
They are uniform over the allowed $q,d$ and page.
The first term in (CW.7) has exponent $802+16\varepsilon$;
the second has exponent $800+41\varepsilon/5$, which is smaller;
the third is smaller as well.  This proves (CW.8).
Only the already established polynomial comparison scale is used.

\proofstep{4. The full common-window quantifiers.}
Fix one finite datum, one $0<\varepsilon<1/256$, and one page as
in AC-K35.  Suppose the replacement gate supplies
$J\geq1$, $X_0$ and $\sigma\in(0,1)$ such that for every
$X\geq X_0$ and every prescribed
$\operatorname{WinPair}(\varepsilon,X^{1/J},q,d)$ there is a single
$W$, $X^{1/J}\leq W$, $W^2\leq X$, with
\[
 I_{399}^{\circ}(H_u,u)\leq(1-\sigma)\exp(u^\varepsilon/2)
                       \quad(W\leq u\leq W^2).
 \tag{CW.13}
\]
Choose $U_0\geq2$ so that the uniform error (CW.8) is at most
$(\sigma/2)\exp(u^\varepsilon/2)$ for every $u\geq U_0$.
This choice depends only on the same fixed datum, $\varepsilon$
and $\sigma$, not on $q,d,u,e,\ell,\eta$ or a separately chosen
window.  Replace $X_0$ by $\max(X_0,U_0^J)$.  Then every such
$W$ is at least $U_0$, and (CW.7), (CW.13) imply
\[
 I_{399}(H_u,u)\leq(1-\sigma/2)\exp(u^\varepsilon/2)
                       \quad(W\leq u\leq W^2).
 \tag{CW.14}
\]
The same $J$, pair, page, host and common $W$ have been used.
If the retained application requires several simultaneous fields
or finite labels on one window, take the maximum of their finitely
many uniform error constants; the same argument uses that same
$W$ and does not intersect independently chosen windows.

Conversely $I_{399}^{\circ}\leq I_{399}$ by nonnegativity, so the
old gate implies (CW.13) with its original slack and parameters.
The analogous little-o statements are equivalent by (CW.8).
Finally apply AC-K35 to (CW.14).  Its actual source integer $R.t$,
all coefficient primes outside the removed labels, ordinary and
mixed local conditions, family transport, and analytic hypotheses
are unchanged.  This proves the stated conditional implication.
\end{proof}

\paragraph{Scope of the reduction.}
This statement removes four actual finite characters at one coefficient
prime and all their positive weights congruent to $267$ modulo $268$.
It does not remove arbitrary finite $p$-conductor twists or continuous
principal-unit characters: such characters have not been identified
with characters of $G_u$.  Indeed $|G_u|=144\cdot2^{h(u)}$ is prime
to $269$, so a $269$-power character factoring through it is trivial.
It does not remove the complementary weights, the other actual finite
characters, or any other coefficient prime.  The mixed realization
$\mathcal R_{\rm cub}$ remains unproved.  The four-character arithmetic
removal uses precisely the proved original strict theorem AC-G142.

\paragraph{Source-reading scope.}
The exact applied dependencies are the complete AC-K25 coefficient and
strict-kernel comparison, AC-K26-local (including its termwise (FL.12)
and absolute summed (FL.4) bounds), AC-K26 (EU.1--EU.2),
AC-K26-elliptic, AC-K27, and the full quantifier and cutoff statement
of AC-K35.  The actual AC-G142 character and original-coefficient statement
is used only at $r=0$, $\kappa=1$; its global full-local strict group
contains the inflated group in (CW.3).  This argument uses those established theorems; it does not
claim a new proof of Poitou--Tate duality or of the imaginary-quadratic
main conjecture, and it does not execute Lean.
